% META: {"id":"1SPE-SUITES-EX-008","chapitre":"1SPE-SUITES","type_objet":"exercice","capacites":["1SPE-SUITES-C2"],"capacites_codes":["C2"],"methodes":["M2"],"parcours":1,"competences":["calculer"],"duree_min":10,"mode_creation":"ex_nihilo","sources_inspiration":[],"parametres_sympy":{"u5":23,"r":4,"p":5},"fichier_tex":"chapitres/1SPE-SUITES/exercices/1SPE-SUITES-EX-008.tex","corrige_tex":"chapitres/1SPE-SUITES/corriges/1SPE-SUITES-CO-008.tex","status":"approved","origin":"generated","status_history":["generated","audited","accepted","approved"],"release_acceptance":"RELEASE_OWNER_BATCH_ACCEPTANCE","acceptance_closure_digest":"sha256:9b3ccf9a81c5520fbb7e03b7b2d3e7bf2057a02834b6d908bf3be5f49f3b553f","acceptance_receipt_digest":"sha256:4fad86f0282e0fa4e0419cca1ad6379ae7af5a280c04a4a90880f90a1c044242"}
% BEGIN-VERIFY
% from sympy import *
% # u_5 = 23, r = 4
% # u_0 = u_5 - 5*r = 23 - 20 = 3
% u5 = 23
% r = 4
% u0 = u5 - 5*r
% assert u0 == 3
% n = symbols('n', integer=True, nonnegative=True)
% u = u0 + n*r
% assert u.subs(n, 5) == 23
% assert u.subs(n, 8) == 35
% END-VERIFY
\begin{exercice}{1SPE-SUITES-EX-008}{1}{10}
Mehdi sait que la suite arithmétique $(u_n)$ a pour raison $r = 4$ et que $u_5 = 23$.
\begin{enumerate}
  \item En utilisant la relation $u_n = u_p + (n-p) \times r$, calculer $u_0$.
  \item Écrire la formule du terme général $u_n$.
  \item Calculer $u_8$.
\end{enumerate}
\end{exercice}
